% =====================================================================
% PREÁMBULO COMÚN — Documentación de Auditoría
% Observatorio MinCiencias · Universidad Santo Tomás · 2026-I
%
% Uso: % =====================================================================
% PREÁMBULO COMÚN — Documentación de Auditoría
% Observatorio MinCiencias · Universidad Santo Tomás · 2026-I
%
% Uso: % =====================================================================
% PREÁMBULO COMÚN — Documentación de Auditoría
% Observatorio MinCiencias · Universidad Santo Tomás · 2026-I
%
% Uso: % =====================================================================
% PREÁMBULO COMÚN — Documentación de Auditoría
% Observatorio MinCiencias · Universidad Santo Tomás · 2026-I
%
% Uso: \input{preambulo.tex} al inicio del preámbulo de cada main.tex
% Compatible con Overleaf (TeX Live completo) y TinyTeX local.
% Paquetes necesarios (instalados vía `tlmgr install` en TinyTeX):
%   tcolorbox pgf environ trimspaces fp caption colortbl enumitem
%   epigraph fancyhdr listings microtype multirow parskip pdflscape
%   siunitx titlesec
% =====================================================================

\documentclass[11pt,a4paper,oneside]{article}

% ---- Codificación e idioma -----------------------------------------
\usepackage[utf8]{inputenc}
\usepackage[T1]{fontenc}
\usepackage[spanish,es-tabla,es-lcroman]{babel}
\usepackage{lmodern}
\usepackage{microtype}

% ---- Símbolos Unicode de respaldo (por si algún texto los contiene) -
\DeclareUnicodeCharacter{2248}{$\approx$}
\DeclareUnicodeCharacter{2265}{$\geq$}
\DeclareUnicodeCharacter{2264}{$\leq$}
\DeclareUnicodeCharacter{2260}{$\neq$}
\DeclareUnicodeCharacter{00B1}{$\pm$}
\DeclareUnicodeCharacter{00D7}{$\times$}
\DeclareUnicodeCharacter{2192}{$\rightarrow$}
\DeclareUnicodeCharacter{2190}{$\leftarrow$}
\DeclareUnicodeCharacter{2194}{$\leftrightarrow$}
\DeclareUnicodeCharacter{21D2}{$\Rightarrow$}
\DeclareUnicodeCharacter{00B7}{$\cdot$}
\DeclareUnicodeCharacter{2026}{\dots{}}
\DeclareUnicodeCharacter{2014}{---}
\DeclareUnicodeCharacter{2013}{--}
\DeclareUnicodeCharacter{2208}{$\in$}
\DeclareUnicodeCharacter{2209}{$\notin$}
\DeclareUnicodeCharacter{222A}{$\cup$}
\DeclareUnicodeCharacter{2229}{$\cap$}
\DeclareUnicodeCharacter{221A}{$\sqrt{}$}
\DeclareUnicodeCharacter{2211}{$\sum$}
\DeclareUnicodeCharacter{222B}{$\int$}
\DeclareUnicodeCharacter{221E}{$\infty$}
\DeclareUnicodeCharacter{2205}{$\emptyset$}
\DeclareUnicodeCharacter{00B0}{$^\circ$}
\DeclareUnicodeCharacter{2261}{$\equiv$}
\DeclareUnicodeCharacter{223C}{$\sim$}
\DeclareUnicodeCharacter{21C6}{$\rightleftarrows$}
\DeclareUnicodeCharacter{25B8}{$\triangleright$}
\DeclareUnicodeCharacter{2713}{$\checkmark$}
\DeclareUnicodeCharacter{2717}{$\times$}
\DeclareUnicodeCharacter{2018}{`}
\DeclareUnicodeCharacter{2019}{'}
\DeclareUnicodeCharacter{201C}{``}
\DeclareUnicodeCharacter{201D}{''}

% ---- Geometría y tipografía ----------------------------------------
\usepackage[a4paper,top=2.6cm,bottom=2.6cm,left=2.7cm,right=2.7cm]{geometry}
\usepackage{setspace}
\setstretch{1.18}
\usepackage{parskip}

% ---- Control de desbordes de línea ---------------------------------
% Evita que líneas o tokens largos (rutas, nombres de archivo) se salgan
% del margen o se solapen con el texto vecino.
\sloppy
\emergencystretch=3em
\hyphenpenalty=1000
\tolerance=2000
% Redefine el guion bajo escapado como punto de quiebre (rutas con _).
\makeatletter
\def\_{\leavevmode \kern.06em\vbox{\hrule width.3em}\allowbreak}
\makeatother

% ---- Colores y enlaces ---------------------------------------------
\usepackage{xcolor}
\definecolor{usta}{RGB}{30,55,128}
\definecolor{accent}{RGB}{183,135,30}
\definecolor{rule}{RGB}{170,170,170}
\definecolor{soft}{RGB}{245,245,247}
\definecolor{bueno}{RGB}{34,139,34}
\definecolor{malo}{RGB}{178,34,34}
\usepackage[colorlinks=true,linkcolor=usta,citecolor=usta,urlcolor=usta]{hyperref}

% ---- Figuras y tablas ----------------------------------------------
\usepackage{graphicx}
\usepackage{caption}
\captionsetup{font=small,labelfont=bf,labelsep=period,justification=raggedright,singlelinecheck=false}
\usepackage{booktabs}
\usepackage{longtable}
\usepackage{array}
\usepackage{tabularx}
\usepackage{multirow}
\usepackage{colortbl}
\usepackage{float}
\usepackage{pdflscape}

% ---- Listas --------------------------------------------------------
\usepackage{enumitem}
\setlist{itemsep=2pt,topsep=4pt,parsep=0pt}
\usepackage{epigraph}

% ---- Encabezados de sección ---------------------------------------
\usepackage{titlesec}
\titleformat{\section}
  {\Large\bfseries\color{usta}}{\thesection}{0.7em}{}[\vspace{-0.2em}\color{rule}\hrule]
\titleformat{\subsection}
  {\large\bfseries}{\thesubsection}{0.6em}{}
\titleformat{\subsubsection}
  {\normalsize\bfseries\color{usta!85}}{\thesubsubsection}{0.5em}{}
\titlespacing*{\section}{0pt}{1.6em}{0.6em}
\titlespacing*{\subsection}{0pt}{1.0em}{0.4em}

% ---- Encabezado y pie ----------------------------------------------
\usepackage{fancyhdr}
\pagestyle{fancy}
\renewcommand{\headrulewidth}{0pt}
\fancyfoot[L]{\footnotesize\color{rule} Documentación de Auditoría · Observatorio MinCiencias · USTA · 2026-I}
\fancyfoot[R]{\footnotesize\color{rule}\thepage}

% ---- Recuadros -----------------------------------------------------
\usepackage[most]{tcolorbox}
\tcbset{
  hallazgo/.style={
    enhanced, colback=soft, colframe=accent, boxrule=0.6pt, arc=2pt,
    left=10pt, right=10pt, top=8pt, bottom=8pt, fontupper=\small,
    title={\color{usta}\bfseries Hallazgo}, coltitle=black,
    breakable
  },
  riesgo/.style={
    enhanced, colback=soft, colframe=malo!60, boxrule=0.6pt, arc=2pt,
    left=10pt, right=10pt, top=8pt, bottom=8pt, fontupper=\small,
    title={\color{malo}\bfseries Riesgo / Limitación}, coltitle=black,
    breakable
  },
  bueno/.style={
    enhanced, colback=soft, colframe=bueno!60, boxrule=0.6pt, arc=2pt,
    left=10pt, right=10pt, top=8pt, bottom=8pt, fontupper=\small,
    title={\color{bueno}\bfseries Fortaleza}, coltitle=black,
    breakable
  },
  nota/.style={
    enhanced, colback=soft, colframe=usta!40, boxrule=0.6pt, arc=2pt,
    left=10pt, right=10pt, top=8pt, bottom=8pt, fontupper=\small,
    title={\color{usta}\bfseries Nota metodológica}, coltitle=black,
    breakable
  }
}

% ---- Código --------------------------------------------------------
\usepackage{listings}
\lstdefinestyle{auditoria}{
  language=Python, basicstyle=\ttfamily\small,
  commentstyle=\color{usta!70}, keywordstyle=\color{usta}\bfseries,
  stringstyle=\color{accent}, numbers=left, numberstyle=\tiny\color{rule},
  frame=single, rulecolor=\color{rule}, breaklines=true,
  showstringspaces=false, tabsize=4, captionpos=b
}
\lstset{style=auditoria}

% ---- Tablas de datos: tipo consola / resumen -----------------------
\newcommand{\celdaok}[1]{\textcolor{bueno}{\textbf{#1}}}
\newcommand{\celdanok}[1]{\textcolor{malo}{\textbf{#1}}}

% ---- Comando para metadatos del documento --------------------------
\newcommand{\docversion}{1.0}
\newcommand{\docfecha}{\today}
\newcommand{\doctitulo}{Documento de Auditoría}
\newcommand{\docsubtitulo}{}
\newcommand{\docautor}{Consultor Estadístico Senior · Auditoría de Proyectos de Datos}

% ---- Portada simple ------------------------------------------------
\newcommand{\hacerportada}{
\begin{titlepage}
\centering
\vspace*{2cm}
{\color{usta}\rule{\linewidth}{1.5pt}}\\[0.6cm]
{\LARGE\bfseries\color{usta} \doctitulo}\\[0.4cm]
{\large \docsubtitulo}\\[1.2cm]
{\color{usta}\rule{\linewidth}{0.5pt}}\\[1.6cm]
\begin{minipage}{0.8\linewidth}\centering
{\normalsize \docautor}\\[0.3cm]
{\small Universidad Santo Tomás · Ustadistica · 2026-I}\\[0.3cm]
{\small Repositorio auditado: \texttt{Observatorio\_Ministerio\_de\_Ciencias\_Grupo8}}\\[0.3cm]
{\small Versión \docversion{} · \docfecha}\\
\end{minipage}\\[1.6cm]
{\color{usta}\rule{\linewidth}{1.5pt}}
\vfill
{\small Documento generado de forma reproducible a partir de la auditoría del repositorio.
La construcción de este documento está documentada en \texttt{documentacion\_auditoria}.}
\end{titlepage}
}
 al inicio del preámbulo de cada main.tex
% Compatible con Overleaf (TeX Live completo) y TinyTeX local.
% Paquetes necesarios (instalados vía `tlmgr install` en TinyTeX):
%   tcolorbox pgf environ trimspaces fp caption colortbl enumitem
%   epigraph fancyhdr listings microtype multirow parskip pdflscape
%   siunitx titlesec
% =====================================================================

\documentclass[11pt,a4paper,oneside]{article}

% ---- Codificación e idioma -----------------------------------------
\usepackage[utf8]{inputenc}
\usepackage[T1]{fontenc}
\usepackage[spanish,es-tabla,es-lcroman]{babel}
\usepackage{lmodern}
\usepackage{microtype}

% ---- Símbolos Unicode de respaldo (por si algún texto los contiene) -
\DeclareUnicodeCharacter{2248}{$\approx$}
\DeclareUnicodeCharacter{2265}{$\geq$}
\DeclareUnicodeCharacter{2264}{$\leq$}
\DeclareUnicodeCharacter{2260}{$\neq$}
\DeclareUnicodeCharacter{00B1}{$\pm$}
\DeclareUnicodeCharacter{00D7}{$\times$}
\DeclareUnicodeCharacter{2192}{$\rightarrow$}
\DeclareUnicodeCharacter{2190}{$\leftarrow$}
\DeclareUnicodeCharacter{2194}{$\leftrightarrow$}
\DeclareUnicodeCharacter{21D2}{$\Rightarrow$}
\DeclareUnicodeCharacter{00B7}{$\cdot$}
\DeclareUnicodeCharacter{2026}{\dots{}}
\DeclareUnicodeCharacter{2014}{---}
\DeclareUnicodeCharacter{2013}{--}
\DeclareUnicodeCharacter{2208}{$\in$}
\DeclareUnicodeCharacter{2209}{$\notin$}
\DeclareUnicodeCharacter{222A}{$\cup$}
\DeclareUnicodeCharacter{2229}{$\cap$}
\DeclareUnicodeCharacter{221A}{$\sqrt{}$}
\DeclareUnicodeCharacter{2211}{$\sum$}
\DeclareUnicodeCharacter{222B}{$\int$}
\DeclareUnicodeCharacter{221E}{$\infty$}
\DeclareUnicodeCharacter{2205}{$\emptyset$}
\DeclareUnicodeCharacter{00B0}{$^\circ$}
\DeclareUnicodeCharacter{2261}{$\equiv$}
\DeclareUnicodeCharacter{223C}{$\sim$}
\DeclareUnicodeCharacter{21C6}{$\rightleftarrows$}
\DeclareUnicodeCharacter{25B8}{$\triangleright$}
\DeclareUnicodeCharacter{2713}{$\checkmark$}
\DeclareUnicodeCharacter{2717}{$\times$}
\DeclareUnicodeCharacter{2018}{`}
\DeclareUnicodeCharacter{2019}{'}
\DeclareUnicodeCharacter{201C}{``}
\DeclareUnicodeCharacter{201D}{''}

% ---- Geometría y tipografía ----------------------------------------
\usepackage[a4paper,top=2.6cm,bottom=2.6cm,left=2.7cm,right=2.7cm]{geometry}
\usepackage{setspace}
\setstretch{1.18}
\usepackage{parskip}

% ---- Control de desbordes de línea ---------------------------------
% Evita que líneas o tokens largos (rutas, nombres de archivo) se salgan
% del margen o se solapen con el texto vecino.
\sloppy
\emergencystretch=3em
\hyphenpenalty=1000
\tolerance=2000
% Redefine el guion bajo escapado como punto de quiebre (rutas con _).
\makeatletter
\def\_{\leavevmode \kern.06em\vbox{\hrule width.3em}\allowbreak}
\makeatother

% ---- Colores y enlaces ---------------------------------------------
\usepackage{xcolor}
\definecolor{usta}{RGB}{30,55,128}
\definecolor{accent}{RGB}{183,135,30}
\definecolor{rule}{RGB}{170,170,170}
\definecolor{soft}{RGB}{245,245,247}
\definecolor{bueno}{RGB}{34,139,34}
\definecolor{malo}{RGB}{178,34,34}
\usepackage[colorlinks=true,linkcolor=usta,citecolor=usta,urlcolor=usta]{hyperref}

% ---- Figuras y tablas ----------------------------------------------
\usepackage{graphicx}
\usepackage{caption}
\captionsetup{font=small,labelfont=bf,labelsep=period,justification=raggedright,singlelinecheck=false}
\usepackage{booktabs}
\usepackage{longtable}
\usepackage{array}
\usepackage{tabularx}
\usepackage{multirow}
\usepackage{colortbl}
\usepackage{float}
\usepackage{pdflscape}

% ---- Listas --------------------------------------------------------
\usepackage{enumitem}
\setlist{itemsep=2pt,topsep=4pt,parsep=0pt}
\usepackage{epigraph}

% ---- Encabezados de sección ---------------------------------------
\usepackage{titlesec}
\titleformat{\section}
  {\Large\bfseries\color{usta}}{\thesection}{0.7em}{}[\vspace{-0.2em}\color{rule}\hrule]
\titleformat{\subsection}
  {\large\bfseries}{\thesubsection}{0.6em}{}
\titleformat{\subsubsection}
  {\normalsize\bfseries\color{usta!85}}{\thesubsubsection}{0.5em}{}
\titlespacing*{\section}{0pt}{1.6em}{0.6em}
\titlespacing*{\subsection}{0pt}{1.0em}{0.4em}

% ---- Encabezado y pie ----------------------------------------------
\usepackage{fancyhdr}
\pagestyle{fancy}
\renewcommand{\headrulewidth}{0pt}
\fancyfoot[L]{\footnotesize\color{rule} Documentación de Auditoría · Observatorio MinCiencias · USTA · 2026-I}
\fancyfoot[R]{\footnotesize\color{rule}\thepage}

% ---- Recuadros -----------------------------------------------------
\usepackage[most]{tcolorbox}
\tcbset{
  hallazgo/.style={
    enhanced, colback=soft, colframe=accent, boxrule=0.6pt, arc=2pt,
    left=10pt, right=10pt, top=8pt, bottom=8pt, fontupper=\small,
    title={\color{usta}\bfseries Hallazgo}, coltitle=black,
    breakable
  },
  riesgo/.style={
    enhanced, colback=soft, colframe=malo!60, boxrule=0.6pt, arc=2pt,
    left=10pt, right=10pt, top=8pt, bottom=8pt, fontupper=\small,
    title={\color{malo}\bfseries Riesgo / Limitación}, coltitle=black,
    breakable
  },
  bueno/.style={
    enhanced, colback=soft, colframe=bueno!60, boxrule=0.6pt, arc=2pt,
    left=10pt, right=10pt, top=8pt, bottom=8pt, fontupper=\small,
    title={\color{bueno}\bfseries Fortaleza}, coltitle=black,
    breakable
  },
  nota/.style={
    enhanced, colback=soft, colframe=usta!40, boxrule=0.6pt, arc=2pt,
    left=10pt, right=10pt, top=8pt, bottom=8pt, fontupper=\small,
    title={\color{usta}\bfseries Nota metodológica}, coltitle=black,
    breakable
  }
}

% ---- Código --------------------------------------------------------
\usepackage{listings}
\lstdefinestyle{auditoria}{
  language=Python, basicstyle=\ttfamily\small,
  commentstyle=\color{usta!70}, keywordstyle=\color{usta}\bfseries,
  stringstyle=\color{accent}, numbers=left, numberstyle=\tiny\color{rule},
  frame=single, rulecolor=\color{rule}, breaklines=true,
  showstringspaces=false, tabsize=4, captionpos=b
}
\lstset{style=auditoria}

% ---- Tablas de datos: tipo consola / resumen -----------------------
\newcommand{\celdaok}[1]{\textcolor{bueno}{\textbf{#1}}}
\newcommand{\celdanok}[1]{\textcolor{malo}{\textbf{#1}}}

% ---- Comando para metadatos del documento --------------------------
\newcommand{\docversion}{1.0}
\newcommand{\docfecha}{\today}
\newcommand{\doctitulo}{Documento de Auditoría}
\newcommand{\docsubtitulo}{}
\newcommand{\docautor}{Consultor Estadístico Senior · Auditoría de Proyectos de Datos}

% ---- Portada simple ------------------------------------------------
\newcommand{\hacerportada}{
\begin{titlepage}
\centering
\vspace*{2cm}
{\color{usta}\rule{\linewidth}{1.5pt}}\\[0.6cm]
{\LARGE\bfseries\color{usta} \doctitulo}\\[0.4cm]
{\large \docsubtitulo}\\[1.2cm]
{\color{usta}\rule{\linewidth}{0.5pt}}\\[1.6cm]
\begin{minipage}{0.8\linewidth}\centering
{\normalsize \docautor}\\[0.3cm]
{\small Universidad Santo Tomás · Ustadistica · 2026-I}\\[0.3cm]
{\small Repositorio auditado: \texttt{Observatorio\_Ministerio\_de\_Ciencias\_Grupo8}}\\[0.3cm]
{\small Versión \docversion{} · \docfecha}\\
\end{minipage}\\[1.6cm]
{\color{usta}\rule{\linewidth}{1.5pt}}
\vfill
{\small Documento generado de forma reproducible a partir de la auditoría del repositorio.
La construcción de este documento está documentada en \texttt{documentacion\_auditoria}.}
\end{titlepage}
}
 al inicio del preámbulo de cada main.tex
% Compatible con Overleaf (TeX Live completo) y TinyTeX local.
% Paquetes necesarios (instalados vía `tlmgr install` en TinyTeX):
%   tcolorbox pgf environ trimspaces fp caption colortbl enumitem
%   epigraph fancyhdr listings microtype multirow parskip pdflscape
%   siunitx titlesec
% =====================================================================

\documentclass[11pt,a4paper,oneside]{article}

% ---- Codificación e idioma -----------------------------------------
\usepackage[utf8]{inputenc}
\usepackage[T1]{fontenc}
\usepackage[spanish,es-tabla,es-lcroman]{babel}
\usepackage{lmodern}
\usepackage{microtype}

% ---- Símbolos Unicode de respaldo (por si algún texto los contiene) -
\DeclareUnicodeCharacter{2248}{$\approx$}
\DeclareUnicodeCharacter{2265}{$\geq$}
\DeclareUnicodeCharacter{2264}{$\leq$}
\DeclareUnicodeCharacter{2260}{$\neq$}
\DeclareUnicodeCharacter{00B1}{$\pm$}
\DeclareUnicodeCharacter{00D7}{$\times$}
\DeclareUnicodeCharacter{2192}{$\rightarrow$}
\DeclareUnicodeCharacter{2190}{$\leftarrow$}
\DeclareUnicodeCharacter{2194}{$\leftrightarrow$}
\DeclareUnicodeCharacter{21D2}{$\Rightarrow$}
\DeclareUnicodeCharacter{00B7}{$\cdot$}
\DeclareUnicodeCharacter{2026}{\dots{}}
\DeclareUnicodeCharacter{2014}{---}
\DeclareUnicodeCharacter{2013}{--}
\DeclareUnicodeCharacter{2208}{$\in$}
\DeclareUnicodeCharacter{2209}{$\notin$}
\DeclareUnicodeCharacter{222A}{$\cup$}
\DeclareUnicodeCharacter{2229}{$\cap$}
\DeclareUnicodeCharacter{221A}{$\sqrt{}$}
\DeclareUnicodeCharacter{2211}{$\sum$}
\DeclareUnicodeCharacter{222B}{$\int$}
\DeclareUnicodeCharacter{221E}{$\infty$}
\DeclareUnicodeCharacter{2205}{$\emptyset$}
\DeclareUnicodeCharacter{00B0}{$^\circ$}
\DeclareUnicodeCharacter{2261}{$\equiv$}
\DeclareUnicodeCharacter{223C}{$\sim$}
\DeclareUnicodeCharacter{21C6}{$\rightleftarrows$}
\DeclareUnicodeCharacter{25B8}{$\triangleright$}
\DeclareUnicodeCharacter{2713}{$\checkmark$}
\DeclareUnicodeCharacter{2717}{$\times$}
\DeclareUnicodeCharacter{2018}{`}
\DeclareUnicodeCharacter{2019}{'}
\DeclareUnicodeCharacter{201C}{``}
\DeclareUnicodeCharacter{201D}{''}

% ---- Geometría y tipografía ----------------------------------------
\usepackage[a4paper,top=2.6cm,bottom=2.6cm,left=2.7cm,right=2.7cm]{geometry}
\usepackage{setspace}
\setstretch{1.18}
\usepackage{parskip}

% ---- Control de desbordes de línea ---------------------------------
% Evita que líneas o tokens largos (rutas, nombres de archivo) se salgan
% del margen o se solapen con el texto vecino.
\sloppy
\emergencystretch=3em
\hyphenpenalty=1000
\tolerance=2000
% Redefine el guion bajo escapado como punto de quiebre (rutas con _).
\makeatletter
\def\_{\leavevmode \kern.06em\vbox{\hrule width.3em}\allowbreak}
\makeatother

% ---- Colores y enlaces ---------------------------------------------
\usepackage{xcolor}
\definecolor{usta}{RGB}{30,55,128}
\definecolor{accent}{RGB}{183,135,30}
\definecolor{rule}{RGB}{170,170,170}
\definecolor{soft}{RGB}{245,245,247}
\definecolor{bueno}{RGB}{34,139,34}
\definecolor{malo}{RGB}{178,34,34}
\usepackage[colorlinks=true,linkcolor=usta,citecolor=usta,urlcolor=usta]{hyperref}

% ---- Figuras y tablas ----------------------------------------------
\usepackage{graphicx}
\usepackage{caption}
\captionsetup{font=small,labelfont=bf,labelsep=period,justification=raggedright,singlelinecheck=false}
\usepackage{booktabs}
\usepackage{longtable}
\usepackage{array}
\usepackage{tabularx}
\usepackage{multirow}
\usepackage{colortbl}
\usepackage{float}
\usepackage{pdflscape}

% ---- Listas --------------------------------------------------------
\usepackage{enumitem}
\setlist{itemsep=2pt,topsep=4pt,parsep=0pt}
\usepackage{epigraph}

% ---- Encabezados de sección ---------------------------------------
\usepackage{titlesec}
\titleformat{\section}
  {\Large\bfseries\color{usta}}{\thesection}{0.7em}{}[\vspace{-0.2em}\color{rule}\hrule]
\titleformat{\subsection}
  {\large\bfseries}{\thesubsection}{0.6em}{}
\titleformat{\subsubsection}
  {\normalsize\bfseries\color{usta!85}}{\thesubsubsection}{0.5em}{}
\titlespacing*{\section}{0pt}{1.6em}{0.6em}
\titlespacing*{\subsection}{0pt}{1.0em}{0.4em}

% ---- Encabezado y pie ----------------------------------------------
\usepackage{fancyhdr}
\pagestyle{fancy}
\renewcommand{\headrulewidth}{0pt}
\fancyfoot[L]{\footnotesize\color{rule} Documentación de Auditoría · Observatorio MinCiencias · USTA · 2026-I}
\fancyfoot[R]{\footnotesize\color{rule}\thepage}

% ---- Recuadros -----------------------------------------------------
\usepackage[most]{tcolorbox}
\tcbset{
  hallazgo/.style={
    enhanced, colback=soft, colframe=accent, boxrule=0.6pt, arc=2pt,
    left=10pt, right=10pt, top=8pt, bottom=8pt, fontupper=\small,
    title={\color{usta}\bfseries Hallazgo}, coltitle=black,
    breakable
  },
  riesgo/.style={
    enhanced, colback=soft, colframe=malo!60, boxrule=0.6pt, arc=2pt,
    left=10pt, right=10pt, top=8pt, bottom=8pt, fontupper=\small,
    title={\color{malo}\bfseries Riesgo / Limitación}, coltitle=black,
    breakable
  },
  bueno/.style={
    enhanced, colback=soft, colframe=bueno!60, boxrule=0.6pt, arc=2pt,
    left=10pt, right=10pt, top=8pt, bottom=8pt, fontupper=\small,
    title={\color{bueno}\bfseries Fortaleza}, coltitle=black,
    breakable
  },
  nota/.style={
    enhanced, colback=soft, colframe=usta!40, boxrule=0.6pt, arc=2pt,
    left=10pt, right=10pt, top=8pt, bottom=8pt, fontupper=\small,
    title={\color{usta}\bfseries Nota metodológica}, coltitle=black,
    breakable
  }
}

% ---- Código --------------------------------------------------------
\usepackage{listings}
\lstdefinestyle{auditoria}{
  language=Python, basicstyle=\ttfamily\small,
  commentstyle=\color{usta!70}, keywordstyle=\color{usta}\bfseries,
  stringstyle=\color{accent}, numbers=left, numberstyle=\tiny\color{rule},
  frame=single, rulecolor=\color{rule}, breaklines=true,
  showstringspaces=false, tabsize=4, captionpos=b
}
\lstset{style=auditoria}

% ---- Tablas de datos: tipo consola / resumen -----------------------
\newcommand{\celdaok}[1]{\textcolor{bueno}{\textbf{#1}}}
\newcommand{\celdanok}[1]{\textcolor{malo}{\textbf{#1}}}

% ---- Comando para metadatos del documento --------------------------
\newcommand{\docversion}{1.0}
\newcommand{\docfecha}{\today}
\newcommand{\doctitulo}{Documento de Auditoría}
\newcommand{\docsubtitulo}{}
\newcommand{\docautor}{Consultor Estadístico Senior · Auditoría de Proyectos de Datos}

% ---- Portada simple ------------------------------------------------
\newcommand{\hacerportada}{
\begin{titlepage}
\centering
\vspace*{2cm}
{\color{usta}\rule{\linewidth}{1.5pt}}\\[0.6cm]
{\LARGE\bfseries\color{usta} \doctitulo}\\[0.4cm]
{\large \docsubtitulo}\\[1.2cm]
{\color{usta}\rule{\linewidth}{0.5pt}}\\[1.6cm]
\begin{minipage}{0.8\linewidth}\centering
{\normalsize \docautor}\\[0.3cm]
{\small Universidad Santo Tomás · Ustadistica · 2026-I}\\[0.3cm]
{\small Repositorio auditado: \texttt{Observatorio\_Ministerio\_de\_Ciencias\_Grupo8}}\\[0.3cm]
{\small Versión \docversion{} · \docfecha}\\
\end{minipage}\\[1.6cm]
{\color{usta}\rule{\linewidth}{1.5pt}}
\vfill
{\small Documento generado de forma reproducible a partir de la auditoría del repositorio.
La construcción de este documento está documentada en \texttt{documentacion\_auditoria}.}
\end{titlepage}
}
 al inicio del preámbulo de cada main.tex
% Compatible con Overleaf (TeX Live completo) y TinyTeX local.
% Paquetes necesarios (instalados vía `tlmgr install` en TinyTeX):
%   tcolorbox pgf environ trimspaces fp caption colortbl enumitem
%   epigraph fancyhdr listings microtype multirow parskip pdflscape
%   siunitx titlesec
% =====================================================================

\documentclass[11pt,a4paper,oneside]{article}

% ---- Codificación e idioma -----------------------------------------
\usepackage[utf8]{inputenc}
\usepackage[T1]{fontenc}
\usepackage[spanish,es-tabla,es-lcroman]{babel}
\usepackage{lmodern}
\usepackage{microtype}

% ---- Símbolos Unicode de respaldo (por si algún texto los contiene) -
\DeclareUnicodeCharacter{2248}{$\approx$}
\DeclareUnicodeCharacter{2265}{$\geq$}
\DeclareUnicodeCharacter{2264}{$\leq$}
\DeclareUnicodeCharacter{2260}{$\neq$}
\DeclareUnicodeCharacter{00B1}{$\pm$}
\DeclareUnicodeCharacter{00D7}{$\times$}
\DeclareUnicodeCharacter{2192}{$\rightarrow$}
\DeclareUnicodeCharacter{2190}{$\leftarrow$}
\DeclareUnicodeCharacter{2194}{$\leftrightarrow$}
\DeclareUnicodeCharacter{21D2}{$\Rightarrow$}
\DeclareUnicodeCharacter{00B7}{$\cdot$}
\DeclareUnicodeCharacter{2026}{\dots{}}
\DeclareUnicodeCharacter{2014}{---}
\DeclareUnicodeCharacter{2013}{--}
\DeclareUnicodeCharacter{2208}{$\in$}
\DeclareUnicodeCharacter{2209}{$\notin$}
\DeclareUnicodeCharacter{222A}{$\cup$}
\DeclareUnicodeCharacter{2229}{$\cap$}
\DeclareUnicodeCharacter{221A}{$\sqrt{}$}
\DeclareUnicodeCharacter{2211}{$\sum$}
\DeclareUnicodeCharacter{222B}{$\int$}
\DeclareUnicodeCharacter{221E}{$\infty$}
\DeclareUnicodeCharacter{2205}{$\emptyset$}
\DeclareUnicodeCharacter{00B0}{$^\circ$}
\DeclareUnicodeCharacter{2261}{$\equiv$}
\DeclareUnicodeCharacter{223C}{$\sim$}
\DeclareUnicodeCharacter{21C6}{$\rightleftarrows$}
\DeclareUnicodeCharacter{25B8}{$\triangleright$}
\DeclareUnicodeCharacter{2713}{$\checkmark$}
\DeclareUnicodeCharacter{2717}{$\times$}
\DeclareUnicodeCharacter{2018}{`}
\DeclareUnicodeCharacter{2019}{'}
\DeclareUnicodeCharacter{201C}{``}
\DeclareUnicodeCharacter{201D}{''}

% ---- Geometría y tipografía ----------------------------------------
\usepackage[a4paper,top=2.6cm,bottom=2.6cm,left=2.7cm,right=2.7cm]{geometry}
\usepackage{setspace}
\setstretch{1.18}
\usepackage{parskip}

% ---- Control de desbordes de línea ---------------------------------
% Evita que líneas o tokens largos (rutas, nombres de archivo) se salgan
% del margen o se solapen con el texto vecino.
\sloppy
\emergencystretch=3em
\hyphenpenalty=1000
\tolerance=2000
% Redefine el guion bajo escapado como punto de quiebre (rutas con _).
\makeatletter
\def\_{\leavevmode \kern.06em\vbox{\hrule width.3em}\allowbreak}
\makeatother

% ---- Colores y enlaces ---------------------------------------------
\usepackage{xcolor}
\definecolor{usta}{RGB}{30,55,128}
\definecolor{accent}{RGB}{183,135,30}
\definecolor{rule}{RGB}{170,170,170}
\definecolor{soft}{RGB}{245,245,247}
\definecolor{bueno}{RGB}{34,139,34}
\definecolor{malo}{RGB}{178,34,34}
\usepackage[colorlinks=true,linkcolor=usta,citecolor=usta,urlcolor=usta]{hyperref}

% ---- Figuras y tablas ----------------------------------------------
\usepackage{graphicx}
\usepackage{caption}
\captionsetup{font=small,labelfont=bf,labelsep=period,justification=raggedright,singlelinecheck=false}
\usepackage{booktabs}
\usepackage{longtable}
\usepackage{array}
\usepackage{tabularx}
\usepackage{multirow}
\usepackage{colortbl}
\usepackage{float}
\usepackage{pdflscape}

% ---- Listas --------------------------------------------------------
\usepackage{enumitem}
\setlist{itemsep=2pt,topsep=4pt,parsep=0pt}
\usepackage{epigraph}

% ---- Encabezados de sección ---------------------------------------
\usepackage{titlesec}
\titleformat{\section}
  {\Large\bfseries\color{usta}}{\thesection}{0.7em}{}[\vspace{-0.2em}\color{rule}\hrule]
\titleformat{\subsection}
  {\large\bfseries}{\thesubsection}{0.6em}{}
\titleformat{\subsubsection}
  {\normalsize\bfseries\color{usta!85}}{\thesubsubsection}{0.5em}{}
\titlespacing*{\section}{0pt}{1.6em}{0.6em}
\titlespacing*{\subsection}{0pt}{1.0em}{0.4em}

% ---- Encabezado y pie ----------------------------------------------
\usepackage{fancyhdr}
\pagestyle{fancy}
\renewcommand{\headrulewidth}{0pt}
\fancyfoot[L]{\footnotesize\color{rule} Documentación de Auditoría · Observatorio MinCiencias · USTA · 2026-I}
\fancyfoot[R]{\footnotesize\color{rule}\thepage}

% ---- Recuadros -----------------------------------------------------
\usepackage[most]{tcolorbox}
\tcbset{
  hallazgo/.style={
    enhanced, colback=soft, colframe=accent, boxrule=0.6pt, arc=2pt,
    left=10pt, right=10pt, top=8pt, bottom=8pt, fontupper=\small,
    title={\color{usta}\bfseries Hallazgo}, coltitle=black,
    breakable
  },
  riesgo/.style={
    enhanced, colback=soft, colframe=malo!60, boxrule=0.6pt, arc=2pt,
    left=10pt, right=10pt, top=8pt, bottom=8pt, fontupper=\small,
    title={\color{malo}\bfseries Riesgo / Limitación}, coltitle=black,
    breakable
  },
  bueno/.style={
    enhanced, colback=soft, colframe=bueno!60, boxrule=0.6pt, arc=2pt,
    left=10pt, right=10pt, top=8pt, bottom=8pt, fontupper=\small,
    title={\color{bueno}\bfseries Fortaleza}, coltitle=black,
    breakable
  },
  nota/.style={
    enhanced, colback=soft, colframe=usta!40, boxrule=0.6pt, arc=2pt,
    left=10pt, right=10pt, top=8pt, bottom=8pt, fontupper=\small,
    title={\color{usta}\bfseries Nota metodológica}, coltitle=black,
    breakable
  }
}

% ---- Código --------------------------------------------------------
\usepackage{listings}
\lstdefinestyle{auditoria}{
  language=Python, basicstyle=\ttfamily\small,
  commentstyle=\color{usta!70}, keywordstyle=\color{usta}\bfseries,
  stringstyle=\color{accent}, numbers=left, numberstyle=\tiny\color{rule},
  frame=single, rulecolor=\color{rule}, breaklines=true,
  showstringspaces=false, tabsize=4, captionpos=b
}
\lstset{style=auditoria}

% ---- Tablas de datos: tipo consola / resumen -----------------------
\newcommand{\celdaok}[1]{\textcolor{bueno}{\textbf{#1}}}
\newcommand{\celdanok}[1]{\textcolor{malo}{\textbf{#1}}}

% ---- Comando para metadatos del documento --------------------------
\newcommand{\docversion}{1.0}
\newcommand{\docfecha}{\today}
\newcommand{\doctitulo}{Documento de Auditoría}
\newcommand{\docsubtitulo}{}
\newcommand{\docautor}{Consultor Estadístico Senior · Auditoría de Proyectos de Datos}

% ---- Portada simple ------------------------------------------------
\newcommand{\hacerportada}{
\begin{titlepage}
\centering
\vspace*{2cm}
{\color{usta}\rule{\linewidth}{1.5pt}}\\[0.6cm]
{\LARGE\bfseries\color{usta} \doctitulo}\\[0.4cm]
{\large \docsubtitulo}\\[1.2cm]
{\color{usta}\rule{\linewidth}{0.5pt}}\\[1.6cm]
\begin{minipage}{0.8\linewidth}\centering
{\normalsize \docautor}\\[0.3cm]
{\small Universidad Santo Tomás · Ustadistica · 2026-I}\\[0.3cm]
{\small Repositorio auditado: \texttt{Observatorio\_Ministerio\_de\_Ciencias\_Grupo8}}\\[0.3cm]
{\small Versión \docversion{} · \docfecha}\\
\end{minipage}\\[1.6cm]
{\color{usta}\rule{\linewidth}{1.5pt}}
\vfill
{\small Documento generado de forma reproducible a partir de la auditoría del repositorio.
La construcción de este documento está documentada en \texttt{documentacion\_auditoria}.}
\end{titlepage}
}
